\chapter{Companion Code and Notebooks}
\label{app:companion}

The numerical claims in the chapters are reproduced by the Python package
\texttt{bernoulli-types} (import name \texttt{bernoulli}) and by notebooks that
accompany the chapters. Both live in the \texttt{bernoulli-py/} directory of the
source repository, \url{https://github.com/queelius/bernoulli}. Paths below are
relative to that directory. The package is MIT-licensed research software at an
alpha stage.

\section*{Installation}

The package requires Python 3.11 or later and has no required runtime
dependencies. From \texttt{bernoulli-py/}:
\begin{lstlisting}[language=bash,numbers=none]
pip install -e ".[notebooks]"  # package, Jupyter, numpy, scipy, ...
pip install -e ".[dev]"        # tests, linting, type checking
pytest                         # the test suite
jupyter lab notebooks/
\end{lstlisting}

\section*{Package modules}

\begingroup
\small
\renewcommand{\arraystretch}{1.15}
\begin{longtable}{@{}>{\raggedright\arraybackslash}p{0.36\textwidth}
                   >{\raggedright\arraybackslash}p{0.44\textwidth}
                   >{\raggedright\arraybackslash}p{0.14\textwidth}@{}}
\toprule
Module (under \texttt{src/bernoulli/}) & Contents & Chapters \\
\midrule
\endfirsthead
\toprule
Module (under \texttt{src/bernoulli/}) & Contents & Chapters \\
\midrule
\endhead
\bottomrule
\endfoot
\texttt{atom.py} & \texttt{BernoulliBool}, the noisy bit with rates
  $(\fprate, \fnrate)$, and single-gate rules & 1--4 \\
\texttt{formula.py} & exact error of a Boolean formula over independent noisy
  bits, conditional on each latent word & 3 \\
\texttt{algebra.py} & composition of rate chains & 3, 6 \\
\texttt{measures.py} & PPV, NPV, $F_1$, Youden's $J$ from
  $(\fprate, \fnrate, \pi)$ & 7 \\
\texttt{algorithms/\allowbreak miller\_rabin.py} & the Miller--Rabin witness test as a
  one-sided noisy bit & 3, 11 \\
\texttt{algorithms/\allowbreak noisy\_classifier.py} & a configurable two-sided noisy
  classifier & 1, 7 \\
\texttt{sets/bloom.py} & Bloom filter as an approximate set & 5 \\
\texttt{sets/\allowbreak composition.py} & rate propagation for Boolean combinations of
  approximate sets & 6 \\
\texttt{sets/cascade.py} & tiered Bloom-cascade membership pipeline & 6, 8 \\
\texttt{sets/\allowbreak entropy\_bound.py} & space and accuracy analyses for one-sided
  approximate sets & 8 \\
\texttt{sets/\allowbreak correlated\_bloom.py} & multi-salt Bloom filter for studying
  correlated hash channels & 16 \\
\texttt{maps/\allowbreak retrieval.py}, \texttt{maps/\allowbreak codec.py} & static GF(2) retrieval
  and value codecs that control non-member outputs & 9, 17 \\
\texttt{maps/\allowbreak perfect\_hash.py} & a BBHash-style minimal perfect hash function
  & 15 \\
\texttt{maps/\allowbreak perfect\_filter.py} & perfect hash filter (minimal perfect hash
  plus fingerprints) & 14, 15 \\
\texttt{relations/\allowbreak relation.py} & approximate relations with a relational
  algebra & 12 \\
\texttt{sketches/\allowbreak countmin.py}, \texttt{hyperloglog.py}, \texttt{minhash.py}
  & Count-Min, HyperLogLog and MinHash as estimators & 10 \\
\texttt{harness.py}, \texttt{bench.py} & empirical rate measurement and
  timing utilities & all \\
\end{longtable}
\endgroup

\section*{Notebooks}

Each notebook opens with a pointer to the chapter it accompanies and, where
practical, asserts the chapter's numbers so that a notebook that runs to
completion certifies them. Chapters~4, 13, 17 and~19 are theoretical and have no
notebook.

\begingroup
\small
\renewcommand{\arraystretch}{1.15}
\begin{longtable}{@{}>{\raggedright\arraybackslash}p{0.10\textwidth}
                   >{\raggedright\arraybackslash}p{0.84\textwidth}@{}}
\toprule
Chapter & Notebook (under \texttt{notebooks/}) \\
\midrule
\endfirsthead
\toprule
Chapter & Notebook (under \texttt{notebooks/}) \\
\midrule
\endhead
\bottomrule
\endfoot
1 & \texttt{ch01-noisy-bit.ipynb} \\
2 & \texttt{ch02-binary-channels.ipynb} \\
3 & \texttt{ch03-boolean-algebra-noise.ipynb} \\
5 & \texttt{ch05-approximate-sets.ipynb} \\
6 & \texttt{ch06-set-composition.ipynb} \\
7 & \texttt{ch07-classification-measures.ipynb} \\
8 & \texttt{ch08-information-bounds.ipynb} \\
9 & \texttt{ch09-approximate-functions.ipynb} \\
10 & \texttt{ch10-sketches-estimators.ipynb} \\
11 & \texttt{ch11-randomized-algorithms.ipynb} \\
12 & \texttt{ch12-approximate-relations.ipynb} \\
14 & \texttt{ch14-bhf.ipynb} \\
15 & \texttt{ch15-bphf.ipynb} \\
16 & \texttt{ch16-not-bernoulli.ipynb} \\
18 & \texttt{ch18-encrypted-search.ipynb} \\
\end{longtable}
\endgroup

\section*{Experiments}

The \texttt{experiments/} directory holds sixteen scripted experiments, each
with a results write-up (\texttt{RESULTS\_NN.md}): composition validation (01),
adversarial Miller--Rabin (02), Bloom-cascade frontiers and operational payoff
(03, 08), codec entropy, optimality and leakage (04, 07, 12 to 16), correlated
hash channels and a conditional-independence audit (05, 06), HyperLogLog as an
estimator (09), Bayesian inference over noisy observations (10), and Le~Cam's
bound as a security limit (11). The write-up for experiment~14 carries a
correction note: an earlier version reversed the direction of the codec
conclusion, and the corrected direction is the one stated in
\Cref{ch:encrypted-search}.

\section*{How the numbers were checked}

Three kinds of check stand behind the quantitative statements. Small cases are
verified by exhaustive enumeration with exact rational arithmetic (for example
the latent-assignment counterexample of \Cref{ex:two-rate-counterexample} and
the order-statistic counterexample of \Cref{ch:optimal-bernoulli-sets}).
Asymptotic and closed-form claims are compared with Monte Carlo estimates at
stated seeds and sample sizes. Claims about prior work are checked against the
cited primary sources. A notebook or test that fails is a reason to distrust the
corresponding sentence of the book until it is resolved.
